% =====================================================================
%  config/dados-da-equipe.tex
%  ESTE É O ÚNICO ARQUIVO DE CONFIGURAÇÃO QUE A EQUIPE PRECISA EDITAR.
%  Leva dois minutos. Faça isso antes de escrever qualquer coisa.
% =====================================================================

% --- 1. Qual ciclo de projeto este relatório documenta? ---------------
%     1 = AP1 (sensoriamento e atuação)
%     2 = AP2 (conectividade e conjunto de dados)
%     3 = AP3 (projeto integrador de AIoT)
%     Trocar este número liga e desliga capítulos inteiros do relatório.
\def\CICLO{2}

% --- 2. Mostrar as instruções cinzas ("Como preencher")? --------------
%     Mantenha ligado enquanto escreve.
%     COMENTE a linha abaixo para gerar a versão de entrega.
% \def\MOSTRARNOTAS{}

% --- 3. Identificação do projeto --------------------------------------
\newcommand{\tituloProjeto}{FitoSfera}
\newcommand{\subtituloProjeto}{nó sensor de temperatura, umidade e luz que entrega suas leituras por Bluetooth Low Energy a quem cuida de um pequeno jardim doméstico}
\newcommand{\nomeEquipe}{FitoSfera}
\newcommand{\repositorioEquipe}{https://github.com/lucasvilari/FitoSfera-Topicos}

% --- 4. Integrantes ---------------------------------------------------
%     Uma linha por integrante: \integrante{Nome completo}{Matrícula}{Papel no ciclo}
%     Papéis do Módulo 1: integração, firmware, hardware, documentação.
%     Papéis do Módulo 2: integração, firmware e comunicação, dados
%     (gravador e qualidade), documentação. Os papéis rotacionam:
%     ninguém repete na AP2 o papel que teve na AP1.
\newcommand{\listaIntegrantes}{%
  %   Na AP1: Lucas = hardware; Gabriel = firmware; Matheus = integração e documentação.
  \integrante{Lucas Vilarinho}{20249011769}{\emph{firmware} e comunicação}
  \integrante{Gabriel Araújo}{20249022010}{documentação}
  \integrante{Matheus Carneiro}{20249012336}{dados e integração}
}

% --- 5. Dados institucionais (normalmente não mudam) ------------------
\newcommand{\instituicaoCurso}{Universidade Federal do Piauí}
\newcommand{\unidadeCurso}{Centro de Ciências da Natureza}
\newcommand{\programaCurso}{Departamento de Computação}
\newcommand{\nomeDisciplina}{Tópicos em Sistemas Computacionais I --- Internet das Coisas e AIoT}
\newcommand{\periodoLetivo}{2026.2}
\newcommand{\docenteDisciplina}{Prof. Atos Apollo Silva Borges}
\newcommand{\localEntrega}{Teresina}
\newcommand{\dataEntrega}{25 de outubro de 2026}

% --- 6. Uso de IA generativa ------------------------------------------
%     A disciplina exige declaração, não proíbe o uso.
%     Preencha o Apêndice D; aqui vai só o resumo de uma linha.
\newcommand{\resumoUsoIA}{%
  A equipe usou IA generativa (Claude, da Anthropic) como apoio no
  código do \emph{firmware}, para escrever o gravador BLE e as
  ferramentas de análise, para operar o gravador e os ensaios
  automáticos na campanha e para redigir este relatório, revisado pela
  equipe. O detalhamento fica no
  Apêndice~\ref{ap:ia}.}
